% Consolidated guide-ch05 from 9a1384ddf1afcd855b09256c9c6f1c90b5023140; source fragments preserved.

% === BEGIN guide-ch05-examples.tex blob 227e4ecee922ef1f3baa388ed1d40a7173e1830f ===
% 《大学物理学学习指导》（第5版）第5章 机械振动，例5.1—5.6。
% 印刷页66—71、PDF页74—79。原教材示例的独立研究型推导；待学术二校。

\sourceexample{5.1}{66}{74}
\problemtext 竖直弹簧劲度系数$k$，原来托住质量$m$的物块$C$；
另一块质量同为$m$的物块$B$从其上方$h$处自由落下，与$C$完全非弹性碰撞。
求两块不分离时的振动表达式及允许的最大落差$h$。
\workings
碰前$B$速度$v_B=\sqrt{2gh}$；碰撞冲量过程两块合速度
$v_0=mv_B/(2m)=\sqrt{gh/2}$（向下）。
以碰后两块压缩弹簧的\emph{新平衡位置}为$x=0$，向下为正，则碰撞瞬间
\[
x_0=-\frac{mg}{k},\qquad
\omega=\sqrt{\frac{k}{2m}},\qquad v_0=\sqrt{\frac{gh}{2}}.
\]
因此不含相位支选择的运动方程为
\[
\boxed{x(t)=-\frac{mg}{k}\cos(\omega t)
       +\frac{\sqrt{gh/2}}{\omega}\sin(\omega t)}
\quad(t\ge0),
\]
振幅
\[
\boxed{A=\sqrt{x_0^2+\frac{v_0^2}{\omega^2}}
=\sqrt{\frac{m^2g^2}{k^2}+\frac{mgh}{k}}}.
\]
为了使上块不脱离下块，二者接触法向力
$N=m(g+\ddot x)$须非负。最大向下加速度为$\omega^2A$，所以
$\omega^2A\le g$，即$A\le2mg/k$。代入得
\[
\boxed{h\le h_{\max}=\frac{3mg}{k}}.
\]
\crosscheck $h=0$时振子起于新平衡位置上方$mg/k$且从静止释放。
若写成余弦相位式，因$x_0<0,v_0>0$，初相必须取第三象限；
只写$\arctan\sqrt{kh/(mg)}$会遗漏$\pi$的象限修正，建议参照本稿无歧义的正余弦分解式。

\sourceexample{5.2}{67}{75}
\problemtext 水平无阻尼弹簧振子质量$M$、劲度系数$k$，振幅$A$；
质量$m$的黏土从上方落到振子上并黏住。分别在振子穿越平衡位置和处于最大位移时碰撞，
求碰后周期和振幅。
\workings
碰前周期$T_0=2\pi\sqrt{M/k}$，黏住后参与振动的质量为$M+m$，
\[
\boxed{T_1=2\pi\sqrt{\frac{M+m}{k}}>T_0}.
\]
若在平衡位置$x=0$碰撞，水平动量守恒使碰后速度
$V=(M/(M+m))v_0$，其中$v_0=\pm A\sqrt{k/M}$。
新角频率$\omega_1=\sqrt{k/(M+m)}$，故新振幅
\[
\boxed{A_{\mathrm{eq}}=\frac{|V|}{\omega_1}
=A\sqrt{\frac{M}{M+m}}}.
\]
若在极端位置碰撞，碰前水平方向速率为零，位移仍为$\pm A$，
所以$\boxed{A_{\mathrm{turn}}=A}$。
\crosscheck 无论在哪个位置，黏连之后周期都变长；水平非弹性碰撞的机械能不能与碰撞前一并视为守恒。

\sourceexample{5.3}{67--68}{75--76}
\problemtext 简谐振动的位移—时间图显示$t=0$时$x=-A/2$且下降，
$t=5\,\mathrm s$时向正方向经过$x=0$。
问振动周期（选项12、10、14、8秒）。
\workings
采用$x=A\cos(\omega t+\varphi)$。
$x(0)=-A/2$且$v(0)<0$意味着$\cos\varphi=-1/2,\sin\varphi>0$，
故$\varphi=2\pi/3$。
$t=5$时向正方向过零，对应第一个与图示连续的相位$3\pi/2$。
\[
5\omega+\frac{2\pi}{3}=\frac{3\pi}{2}
\quad\Longrightarrow\quad
\omega=\frac{\pi}{6}\ \mathrm{rad/s}.
\]
于是
\[
\boxed{T=\frac{2\pi}{\omega}=12\,\mathrm s}\quad(\text{选A}).
\]
\crosscheck 5秒内的相位推进为$5\pi/6$，位移从负半振幅先继续下降到负峰，再回升穿越原点，与原图一致。

\sourceexample{5.4}{68}{76}
\problemtext 质量$100\,\mathrm g$的物体悬挂于竖直弹簧下端，
从静平衡位置向下拉$10\,\mathrm{cm}$后静止释放，弹簧劲度系数$1.0\,\mathrm{N/m}$。
求第一次经过平衡位置的速率。
\workings
以平衡位置为零点，振幅$A=0.10\,\mathrm m$、
运动质量$m=0.10\,\mathrm{kg}$，因此
\[
\omega=\sqrt{\frac{k}{m}}=\sqrt{10}\,\mathrm{rad/s},\quad
\boxed{v_{\max}=A\omega=0.10\sqrt{10}\approx0.316\,\mathrm{m/s}}.
\]
\crosscheck 通过平衡位置时弹簧振动势能为零，振动动能最大；
也可由$kA^2/2=mv_{\max}^2/2$复算。

\sourceexample{5.5}{68--70}{76--78}
\problemtext 已知速度—时间曲线：$t=0$时$v_0=-15.7\,\mathrm{cm/s}$且继续变负，
$t=1\,\mathrm s$时$v_1=+15.7\,\mathrm{cm/s}$且速度随时间增大；
速度振幅$31.4\,\mathrm{cm/s}$。求位移振动方程。
\workings
设$x=A\cos(\omega t+\varphi)$，则
$v=-\omega A\sin(\omega t+\varphi)$。
由$|v|_{\max}=\omega A=31.4$和$t=0$的速度，
$\sin\varphi=1/2$。又因$v$初始变负，$a(0)<0$，故$x(0)>0$，
即$\cos\varphi>0$，从而$\varphi=\pi/6$。
$t=1$时速度为正且正在增大，意味着相位取$7\pi/6$，所以
$\omega+\pi/6=7\pi/6$，得$\omega=\pi\,\mathrm{rad/s}$。
振幅$A=31.4/\pi\,\mathrm{cm}\approx10.0\,\mathrm{cm}$，因此
\[
\boxed{x(t)\approx10\cos\left(\pi t+\frac{\pi}{6}\right)\ \mathrm{cm}}.
\]
\crosscheck $v(0)\approx-15.7$ cm/s；$v(1)\approx+15.7$ cm/s，且$t=0$时加速度为负、$t=1$时加速度为正，与速度曲线局部斜率吻合。

\sourceexample{5.6}{70--71}{78--79}
\problemtext 同方向同频简谐振动
$x_1=6\cos(2t+\pi/6)$ cm，
$x_2=8\cos(2t-\pi/3)$ cm；
用旋转矢量法求合振动振幅与初相。
\workings
两初相差为$(\pi/6)-(-\pi/3)=\pi/2$，旋转矢量互相垂直，所以
\[
\boxed{A=\sqrt{6^2+8^2}=10\,\mathrm{cm}}.
\]
相量直角坐标为
\[
X=6\cos\frac{\pi}{6}+8\cos\left(-\frac{\pi}{3}\right)
=3\sqrt3+4,
\]
\[
Y=6\sin\frac{\pi}{6}+8\sin\left(-\frac{\pi}{3}\right)
=3-4\sqrt3.
\]
故
\[
\boxed{\varphi=\operatorname{atan2}(3-4\sqrt3,\,3\sqrt3+4)
\approx-0.403\,\mathrm{rad}\approx-23.1^\circ},
\]
\[
\boxed{x(t)=10\cos(2t+\varphi)\ \mathrm{cm}}.
\]
\crosscheck $X^2+Y^2=100$；在$t=0$处合位移应为$X\approx9.196$ cm，而非10 cm，符合合振动非零初相。
% === END guide-ch05-examples.tex ===
